\section{Exact third growth on the moving second holding state}
\label{tg:section}

The third growth begins at the actual second return. Its inherited
first vorticity is positive, its second vorticity is zero, and its
second temperature and phase length continue changing. We derive the
complete third phase and linear pair on the holding interval already
proved in Theorem~\ref{sr:theorem}. The original source growth control
is $\dot\alpha=\mathcal F_2$; the different control
$\mathcal F_3$ is retained below as a computed function. This distinction
matters at the next transition. All formulas below use the physical
frequency $K_j=\mu\widehat\lambda_j$, the original
$\nu=\nu_{\rm AB}=\sqrt{A_0}$, and the actual changed scales.

\subsection{Actual birth data and their branches}

Set $t_0=t_2^b$, $\sigma_1=\sigma$ and
$s_2=L_2\nu/\sigma_1$, with $a=\sin s_2$ as in
\eqref{sr:parameters}. Write $h_b=h(t_0)$, $\chi_b=\sqrt{h_b^2+a^2}$,
$P_2=-\Theta_2$ and $A_1=-K_1\Theta_1$. On the proved interval
\begin{equation}\label{tg:interval}
 I=\left[t_0,t_0+\frac1{32\Gamma_2}\right],\qquad
 \Gamma_2=\sigma_1a,
\end{equation}
the actual physical holding equations and phases are
\begin{equation}\label{tg:parent}
\begin{gathered}
 \chi=\sqrt{h^2+a^2},\quad
 \zeta_1=\frac{(-a,h)}\chi,\quad \zeta_2=(0,\chi),
 \quad\theta=\arctan(a/h),\quad\alpha=s_*+\theta,\\
 \dot h=a\Omega_1,\qquad
 \dot\Omega_1=\frac{aA_1}{\chi}-dK_1^2\Omega_1,\qquad
 \dot A_1=-A_0\sin s_*\Omega_1-dK_1^2 A_1,\\
 \Omega_2=0,\qquad
 \dot P_2=-dK_2^2\chi^2P_2,\qquad
 \dot\alpha=-\frac{a^2\Omega_1}{\chi^2}.
\end{gathered}
\end{equation}
Here $A_1,P_2,\Omega_1$ are strictly positive throughout $I$;
$h_b\leq h<3$, $h_b\geq\cos s_2\geq15/16$ and
$1\leq\chi<\sqrt{10}$ are proved bounds, not frozen values.
The notation $\theta$ in this subsection denotes an angle only;
the temperature field is denoted $\vartheta$ below where needed.

Retain the original recursion
\begin{equation}\label{tg:third_design}
 \begin{gathered}
 Q_3=Q_*+3,\quad
 \widehat\lambda_3=\widehat\lambda_2^{Q_3},\quad
 k_3=\max\{k\in\mathbb N_0:120k\leq Q_3,
                  \ K(k)\leq\log\widehat\lambda_3\},\\
 L_3=(k_3+41/8)\log\widehat\lambda_3,\qquad
 \sigma_2^2=|G_{<3}(t_0)|,\qquad
 s_3=\frac{L_3\sigma_1}{\sigma_2},\qquad
 \Gamma_3=\sigma_2\sin s_3,\qquad p_3=e(s_3).
 \end{gathered}
\end{equation}
The positive square root defining $\sigma_2$ is the source scale
evaluated on the actual state. It is not the magnitude of only the
second temperature contribution. The set defining $k_3$ is nonempty
because $k=0$ belongs to it by \eqref{ave:Ychoice}, and finite because
$120k\leq Q_3$.

For completeness, even the earlier frequency choice fixes the small
positive branch needed for the exact map. Indeed $\alpha\leq1/4$
on $I$, so the vertical component of
$G_{<3}=-A_0R_\alpha e_2-A_1\zeta_1-K_2P_2\zeta_2$
at birth has three negative summands. Therefore
\[
 \sigma_2^2\geq K_2P_2(t_0)\chi_b
 \geq A_0\widehat\lambda_2^{1/8},\qquad
 \sigma_2\geq\nu\widehat\lambda_2^{1/16}.
\]
The threshold and selected gain in Theorem~\ref{ave:complete_entry}
justify the second inequality. Put $y=\log\widehat\lambda_1$.
Equations~\eqref{ave:scale_bounds} and \eqref{ave:Ychoice} give
$\sigma_1\leq C_\sigma\nu e^{y/16}$ and $y\geq Q_2\geq202$.
Since $\Lambda_1\geq6$, one has
$C_\sigma=\sqrt{1+e^{2+\Lambda_1^{-1}}}<6$.
One may verify the numerical bound by $e<3$, which follows from
$n!\geq2^{n-1}$ for $n\geq2$ and strict inequality for $n\geq3$.
Also $Q_3=Q_2+1\leq2Q_2$ and
$k_3+41/8\leq Q_3/120+6\leq Q_3$; hence
$L_3\leq8Q_2^3y$. Consequently
\begin{equation}\label{tg:birth_angle}
 0<s_3\leq48Q_2^3y e^{-(Q_2-1)y/16}
 \leq48y^4e^{-12y}
 \leq\frac{5760}{12^5y}<\frac18.
\end{equation}
The penultimate inequality is the fifth positive term in the
exponential series. Thus $0<\theta_b:=\arctan(a/h_b)\leq s_2<1/8$
and $0<s_3+\theta_b<1/4$. These are statements about the original
source angle, including its branch. Stronger angle and physical
diffusion bounds needed to reach the third threshold are supplied
by the accompanying actual third-growth theorem; none is presumed
in the identities below.

Define the oriented constants and the changing inner product by
\begin{equation}\label{tg:constants}
 b=\sin(s_3+\theta_b),\qquad g_b=\cos(s_3+\theta_b),\qquad
 g=g_b+\frac ba(h-h_b),\qquad
 c=bh_b-ag_b=\chi_b\sin s_3>0.
\end{equation}
The constant $b$ in \eqref{tg:constants} is dimensionless and must not
be confused with the physical buoyancy coefficient $\mathfrak b_3$
below. One has $b>0$, $g_b>15/16$, $g\geq g_b$, and the exact identity
$bh-ag=c$ on $I$.

\subsection{The complete material map, phases and original feedback}

For $v\in\mathbb R^2$, the matrix convention is
$Jv=(-v_2,v_1)$ and $R_\beta=\exp(\beta J)$.
The third deformation and its inverse are
\begin{equation}\label{tg:material}
\begin{aligned}
 M_3(t)&=M_2(t)M_2(t_0)^{-1}
 =R_\theta
   \begin{pmatrix}1&-(h-h_b)/a\\0&1\end{pmatrix}R_{-\theta_b},\\
 M_3(t)^{-1}&=R_{\theta_b}
   \begin{pmatrix}1&(h-h_b)/a\\0&1\end{pmatrix}R_{-\theta}.
\end{aligned}
\end{equation}
These formulas describe linear maps $\mathbb R^2\to\mathbb R^2$
in the original physical coordinates and have determinant one.
In particular $M_3(t_0)=I$ retains the source third-birth datum;
it does not reset $M_1$ or $M_2$.
Their phase pullback gives
\begin{equation}\label{tg:phase}
 \begin{gathered}
 \zeta_3=M_3^{-T}p_3=g\zeta_1-bJ\zeta_1
      =\frac{(c,hg+ab)}\chi,\qquad
 R_3:=|\zeta_3|^2=g^2+b^2,\\
 \det(\zeta_1,\zeta_3)=-b,\qquad
 \det(\zeta_2,\zeta_3)=-c,\qquad
 \dot g=b\Omega_1,\quad \dot R_3=2bg\Omega_1,
 \quad\dot\chi=\frac{ah\Omega_1}{\chi}.
 \end{gathered}
\end{equation}
Thus $R_3(t_0)=1$ and $R_3\geq1$. Its first component is positive
and decreases on growth; it is not held constant by the preceding
feedback. The lower component is positive and satisfies the explicit
bound
\begin{equation}\label{tg:domain}
 \zeta_{3,2}=\frac{hg+ab}{\chi}
       >\frac{225}{256\sqrt{10}}>\frac14.
\end{equation}
For example the last strict comparison follows on squaring from
$225^2>64^2\,10$.

\begin{proof}[Proof of the geometric map]
On holding $\Omega_2=0$ and the older cutoffs equal one, so
$D_{<3}=D_{<2}=\dot\alpha J+
\Omega_1J\zeta_1\otimes\zeta_1$.
The previously proved formula for $M_2$ is
$R_\theta\left(\begin{smallmatrix}1&-(h-\cos s_2)/a\\0&1\end{smallmatrix}\right)$.
Multiplication by its birth inverse gives \eqref{tg:material}.
Differentiating $M_2(t)M_2(t_0)^{-1}$ gives
$\dot M_3=D_{<3}M_3$, including the original initial value.
Multiplying the proposed inverse by $M_3$ on either side gives $I$.

Since $R_{-\theta_b}e(s_3)=(b,g_b)$, inverse transposition
of the shear sends this vector to $(b,g)$. Subsequent rotation gives
$R_\theta(b,g)=(-ag+bh,hg+ab)/\chi$, proving the pullback in
\eqref{tg:phase}. Because $\zeta_1=R_\theta e_2$ and
$J\zeta_1=R_\theta(-e_1)$, this is also
$g\zeta_1-bJ\zeta_1$. Its norm, inner product and signed
determinants follow by the displayed orthonormal basis and the two
laboratory components. Differentiating $bh-ag=c$ and using
$\dot h=a\Omega_1$ gives $\dot g=b\Omega_1$.
Direct differentiation gives the remaining rates.
All maps are invertible in their specified domains because
$a,h,\chi>0$; their branch choices were proved above.
The lower bound \eqref{tg:domain} follows from
$h,g>15/16$ and $\chi<\sqrt{10}$.
\end{proof}

Write the original feedback functions at the actual state as
\begin{equation}\label{tg:feedback}
 \mathcal F_2=-\frac{a^2\Omega_1}{\chi^2},\qquad
 \mathcal F_3=-\frac{ab\Omega_1}{hg+ab},\qquad
 \mathcal F_3-\mathcal F_2
    =-\frac{ach\Omega_1}{\chi^2(hg+ab)}.
\end{equation}
To verify them from the original sum, use
\[
 \mathcal F_j=-\frac{\sum_{i<j}\Omega_i e_{i,1}
                         (Je_i\cdot\zeta_j)}{\zeta_{j,2}}.
\]
Here $e_1=\zeta_1$, $J\zeta_1\cdot\zeta_2=-a$,
$J\zeta_1\cdot\zeta_3=-b$ and $\Omega_2=0$.
The two substitutions give the first two formulas in
\eqref{tg:feedback}. Taking their difference and using $bh-ag=c$
gives the third. In particular the physical controlled-component
identity is exactly
\[
 \dot\zeta_{3,1}=(\mathcal F_3-\mathcal F_2)\zeta_{3,2}
             =-\frac{ach\Omega_1}{\chi^3}
             =\frac{d}{dt}\left(\frac c\chi\right).
\]
No term has been discarded by replacing the growth feedback with
the third feedback. Conversely, substituting \eqref{tg:feedback}
in the full phase equation reconstructs this same growth map.

\subsection{Both inherited temperatures and the exact numerator identity}

Put $S_*=A_0\sin s_*$ and $C_*=A_0\cos s_*$. These symbols
retain the original constants; they are not new adjustable scales.
The complete parent gradient can be written in two identical forms:
\begin{equation}\label{tg:gradient}
\begin{aligned}
 G_{<3}&=-A_0R_\alpha e_2-A_1\zeta_1-K_2P_2\zeta_2\\
       &=-(A_1+C_*+K_2P_2h)\zeta_1
                   -(S_*-aK_2P_2)J\zeta_1.
\end{aligned}
\end{equation}
The second form follows from
$R_\alpha e_2=\cos s_*\zeta_1+\sin s_*J\zeta_1$ and
$\zeta_2=h\zeta_1-aJ\zeta_1$. It also supplies the exact
source scale formula
\begin{equation}\label{tg:scale}
 \sigma_2^4=
  (A_1(t_0)+C_*+K_2P_2(t_0)h_b)^2
  +(S_*-aK_2P_2(t_0))^2.
\end{equation}
There is no sign restriction on the second bracket, and none is
needed in this identity.

Define
\begin{equation}\label{tg:numerator}
 N_3=b(A_1+C_*)+gS_*+cK_2P_2,\qquad
 E_3=bK_1^2A_1+cK_2^3\chi^2P_2.
\end{equation}
Every summand of $N_3$ and $E_3$ is positive on $I$.
The exact scalar-temperature and vorticity couplings are
\begin{equation}\label{tg:couplings}
 \mathfrak a_3=-\frac{J\zeta_3\cdot G_{<3}}{K_3R_3}
              =\frac{N_3}{K_3R_3}>0,
 \qquad
 \mathfrak b_3=K_3\zeta_{3,1}=\frac{K_3c}{\chi}>0.
\end{equation}
Their exact rates, including all changing-parent contributions, are
\begin{equation}\label{tg:coefficient_rates}
\begin{gathered}
 \dot N_3=-dE_3,\qquad
 N_3(t)=N_3(t_0)-d\int_{t_0}^t
       \{bK_1^2A_1(r)+cK_2^3\chi(r)^2P_2(r)\}\,dr,\\
 \frac{d}{dt}\log\mathfrak a_3
       =-\frac{dE_3}{N_3}-\frac{2bg\Omega_1}{R_3},\qquad
 \frac{d}{dt}\log\mathfrak b_3
       =-\frac{ah\Omega_1}{\chi^2}.
\end{gathered}
\end{equation}
In particular the numerator would be exactly constant at $d=0$
on this holding geometry. With positive physical diffusivity its
replacement is the integral in \eqref{tg:coefficient_rates}, with
both older temperature losses retained. Its positivity follows from
the original sum in \eqref{tg:numerator}, not from assuming that this
integral is small.

\begin{proof}
Using \eqref{tg:gradient} and
$J\zeta_3=gJ\zeta_1+b\zeta_1$ gives
\[
 -J\zeta_3\cdot G_{<3}
 =b(A_1+C_*+K_2P_2h)+g(S_*-aK_2P_2)=N_3.
\]
This proves the full coupling and its signs. Next differentiate the
three summands of $N_3$, using \eqref{tg:parent} and
\eqref{tg:phase}. The two terms
$-bS_*\Omega_1$ and $\dot gS_*=bS_*\Omega_1$ cancel exactly;
the surviving terms are $-dbK_1^2A_1$ and
$-dcK_2^3\chi^2P_2$. Thus $\dot N_3=-dE_3$.
Integrating proves its displayed integral identity. Finally apply
the logarithmic derivative to \eqref{tg:couplings} and insert
$\dot R_3$, $\dot\chi$ and $\dot N_3$. All logarithms are defined
because their arguments have already been proved positive.
\end{proof}

\subsection{Exact third amplitudes, inverse maps and fundamental functions}

The original signed seed and activation are
\begin{equation}\label{tg:seed}
\begin{gathered}
 S_3=\frac{A_0}{\mu}\widehat\lambda_3^{-k_3-6},\quad
 u_i=\sqrt{\mathfrak b_3(t_0)/\mathfrak a_3(t_0)},\\
 \Theta_3(t_0)=-S_3,\quad \Omega_3(t_0)=-S_3u_i,\quad
 w_3(t)=\eta(\Gamma_3(t-t_0)).
\end{gathered}
\end{equation}
The homogeneous pair is active in the coefficient construction at
$t_0$, while its physical increment is multiplied by $w_3$, exactly
as in source (3.11). The full physical equal-diffusion equations are
\begin{equation}\label{tg:pair}
 \dot\Theta_3=\mathfrak a_3\Omega_3-dK_3^2R_3\Theta_3,
 \qquad
 \dot\Omega_3=\mathfrak b_3\Theta_3-dK_3^2R_3\Omega_3.
\end{equation}

Here are convergent fundamental functions for the changing coefficients,
so the solution does not rely on a constant-coefficient surrogate.
Let
\begin{equation}\label{tg:fundamental}
 \begin{gathered}
 \mathcal A(t)=\begin{pmatrix}0&\mathfrak a_3(t)\\
                              \mathfrak b_3(t)&0\end{pmatrix},\quad
 \mathcal H_0(t)=I,\qquad
 \mathcal H_n(t)=\int_{t_0}^t\mathcal A(r)\mathcal H_{n-1}(r)\,dr,\\
 \mathcal H(t)=\sum_{n=0}^\infty\mathcal H_n(t),\qquad
 \Lambda_3(t)=dK_3^2\int_{t_0}^tR_3(r)\,dr,\qquad
 \mathcal U(t)=e^{-\Lambda_3(t)}\mathcal H(t).
 \end{gathered}
\end{equation}
Each term is an explicit ordered integral of the actual coefficients.
The solution and inverse-data map are
\begin{equation}\label{tg:fundamental_inverse}
\begin{gathered}
 \binom{\Theta_3(t)}{\Omega_3(t)}
       =-S_3\mathcal U(t)\binom1{u_i},\qquad
 \binom{\Theta_3(t_0)}{\Omega_3(t_0)}
       =e^{\Lambda_3(t)}\mathcal H(t)^{-1}
                         \binom{\Theta_3(t)}{\Omega_3(t)},\\
 \det\mathcal U=e^{-2\Lambda_3}>0.
\end{gathered}
\end{equation}

\begin{proof}[Existence, inversion and signs]
All parent coefficients are smooth on the compact interval $I$, and
their denominators have the positive bounds already proved. For any
matrix operator norm put $M=\max_I\|\mathcal A\|$ and
$\ell=|I|$. Repeated integration proves
$\|\mathcal H_n(t)\|\leq M^n(t-t_0)^n/n!$ by induction.
The scalar exponential series therefore proves absolute uniform
convergence. Also
$\mathcal H_n'=\mathcal A\mathcal H_{n-1}$ for $n\geq1$;
the same bound with index $n-1$ proves uniform convergence of the
derivative series. Termwise differentiation gives
$\mathcal H'=\mathcal A\mathcal H$ and $\mathcal H(t_0)=I$.
If two fundamental solutions had the same data, their difference
$V$ would obey $V(t)=\int_{t_0}^t\mathcal A(r)V(r)\,dr$.
Iteration bounds its norm by
$\sup_I\|V\|(M\ell)^n/n!$ for every $n$, which tends to zero;
this proves uniqueness without assuming invertibility first.

For $\mathcal H=\left(\begin{smallmatrix}f_1&f_2\\g_1&g_2\end{smallmatrix}\right)$,
differentiation of $f_1g_2-f_2g_1$ with
$\dot f_j=\mathfrak a_3g_j$,
$\dot g_j=\mathfrak b_3f_j$ gives zero. Its initial value is one.
Thus its inverse is exactly
$\left(\begin{smallmatrix}g_2&-f_2\\-g_1&f_1\end{smallmatrix}\right)$.
Differentiating $\mathcal U$ now proves the full damped matrix
equation; multiplying by $e^{\Lambda_3}$ proves its inverse and
the determinant in \eqref{tg:fundamental_inverse}.
Every entry in $\mathcal A$ is nonnegative, and its two off-diagonal
entries are positive. Consequently every ordered-integral term has
nonnegative entries, the diagonal of $\mathcal H$ is at least one,
and its off-diagonal entries are positive for $t>t_0$ because already
the first term has positive such entries. Since $S_3,u_i>0$, both
amplitudes in \eqref{tg:fundamental_inverse} are strictly negative
throughout $I$. Hence the temperature sign and all following
quotients are established on the complete interval.
\end{proof}

An equivalent scalar fundamental equation retains the geometry in a
particularly direct form. Put $Y=e^{\Lambda_3}\Omega_3$. Then
\begin{equation}\label{tg:scalar_equation}
 \frac d{dt}(\chi\dot Y)-\frac{cN_3}{R_3}Y=0,
 \qquad
 \Omega_3=e^{-\Lambda_3}Y,\qquad
 \Theta_3=e^{-\Lambda_3}\frac{\chi\dot Y}{K_3c}.
\end{equation}
Indeed the first transformed equation is
$\dot X=\mathfrak a_3Y$, where $X=e^{\Lambda_3}\Theta_3$,
and the second is $\dot Y=(K_3c/\chi)X$. The latter gives
$\chi\dot Y=K_3cX$; differentiating it and substituting the former
gives \eqref{tg:scalar_equation}. Conversely that scalar equation
and its displayed inverse give both transformed first-order equations,
and hence \eqref{tg:pair}. In physical damped variables the same
operator is exactly
\begin{equation}\label{tg:damped_scalar}
 (\partial_t+\lambda_3)
       \{\chi(\partial_t+\lambda_3)\Omega_3\}
       -\frac{cN_3}{R_3}\Omega_3=0,
 \qquad \lambda_3=dK_3^2R_3.
\end{equation}
Its complete expanded left side is
\[
 \chi\ddot\Omega_3+(\dot\chi+2\chi\lambda_3)\dot\Omega_3+
 (\chi\dot\lambda_3+\dot\chi\lambda_3+\chi\lambda_3^2-cN_3/R_3)\Omega_3.
\]
In particular $\dot\lambda_3=2dK_3^2bg\Omega_1$ is present in
the physical second-order equation. The integrating factor cancels
the newest diagonal damping on the same evolved parent; it does not
replace $N_3,h,P_2$ by their inviscid counterparts.

\subsection{Exact quotient, gain and diffusion accounting}

Set $\rho=\Gamma_3(t-t_0)$,
$v_3=\sigma_2\Omega_3/(K_3\Theta_3)$ and
$P_3=-\Theta_3$. This change has inverse
$t=t_0+\rho/\Gamma_3$, $\Theta_3=-P_3$,
$\Omega_3=-K_3P_3v_3/\sigma_2$. Define
\begin{equation}\label{tg:quotient_coefficients}
 z_3=\frac{\chi_b}{\chi},\qquad
 k_3^{\rm dyn}=\frac{N_3}{\sigma_2^2\sin s_3\,R_3},\qquad
 \delta_3=\frac{dK_3^2}{\Gamma_3}.
\end{equation}
The symbol $k_3^{\rm dyn}$ is a time-dependent coupling and is
distinct from the original integer $k_3$ in \eqref{tg:third_design}.
Direct quotient and product differentiation prove the exact mutually
inverse equations
\begin{equation}\label{tg:quotient}
 \frac{dv_3}{d\rho}=z_3-k_3^{\rm dyn}v_3^2,\qquad
 \frac{d}{d\rho}\log P_3=k_3^{\rm dyn}v_3-\delta_3R_3,
 \qquad
 v_3(0)=\frac{\sigma_2u_i}{K_3}.
\end{equation}
Thus $z_3$ is the changing preceding phase-length ratio, not a
prescribed constant control. The direct diffusion term cancels from
the quotient only because both physical diffusivities equal $d$.
For unequal physical diffusivities $\kappa$ and $\varepsilon$ in
the newest pair on this same parent, the corresponding exact term
is $(\kappa-\varepsilon)K_3^2R_3v_3/\Gamma_3$, and the logarithmic
loss is $\kappa K_3^2R_3/\Gamma_3$. These are algebraic maps for
that pair; they make no assertion that the older equal-diffusion
trajectory solves the unequal-diffusion parent equations.

The instantaneous eigenline formulation used by the source is also
exact. With
\[
 \gamma_3=\sqrt{\mathfrak a_3\mathfrak b_3}
          =\sqrt{\frac{cN_3}{\chi R_3}},\quad
 u_* =\sqrt{\mathfrak b_3/\mathfrak a_3}
          =K_3\sqrt{\frac{cR_3}{\chi N_3}},\quad
 Z=\frac{\Omega_3}{u_*\Theta_3},
\]
one has, in physical time,
\begin{equation}\label{tg:eigenline}
\begin{gathered}
 \dot Z=\gamma_3(1-Z^2)-Z\frac d{dt}\log u_*,
 \qquad
 \frac d{dt}\log u_*
    =\frac{bg\Omega_1}{R_3}
       -\frac{ah\Omega_1}{2\chi^2}+\frac{dE_3}{2N_3},\\
 \frac d{dt}\log P_3=\gamma_3Z-dK_3^2R_3.
\end{gathered}
\end{equation}
These follow by dividing \eqref{tg:pair}, substituting
\eqref{tg:coefficient_rates}, and differentiating $u_*$; the seed
has exactly $Z(t_0)=1$.

For every $t\in I$ the full new-layer diffusion exponent is
\begin{equation}\label{tg:diffusion}
\begin{gathered}
 \Lambda_3(t)=dK_3^2\int_{t_0}^t
       \left\{\left[g_b+\frac ba(h(r)-h_b)\right]^2+b^2\right\}dr,\\
 dK_3^2(t-t_0)\leq\Lambda_3(t)
        \leq dK_3^2R_{3,\max}(t-t_0),
\end{gathered}
\end{equation}
where the finite explicit bound is
$R_{3,\max}=[g_b+(b/a)(3-h_b)]^2+b^2$.
The inequalities follow from $h_b\leq h<3$ and $b,g_b>0$.
The exact gain is the integral
\[
 \log\frac{P_3(t)}{S_3}
    =\int_{t_0}^t\gamma_3(r)Z(r)\,dr-\Lambda_3(t).
\]
The two inherited diagonal damping exponents on this same interval are
\[
 dK_1^2(t-t_0),\qquad dK_2^2\int_{t_0}^t\chi(r)^2dr.
\]
The first temperature also retains the exact response
\[
 A_1(t)=e^{-dK_1^2(t-t_0)}A_1(t_0)
       -S_*\int_{t_0}^t e^{-dK_1^2(t-r)}\Omega_1(r)\,dr.
\]
These expressions retain all three physical frequencies and the
phase-length factors. The formula for $N_3$ records their combined
effect on the third gain coefficient. No attainment of the third
threshold is inferred from positivity of the cooperative transformed
pair alone; actual attainment uses the quantitative estimates proved
in the accompanying continuation.

\subsection{The exact affine parent equations during third growth}

The moving affine parent also has an exact spatial interpretation
throughout $I$. Put
\[
 D=D_{<3},\quad B_1=-A_1\zeta_1,\quad B_2=-K_2P_2\zeta_2,
 \quad\vartheta(x,t)=G_{<3}(t)\cdot x,\quad u(x,t)=D(t)x.
\]
Then
\begin{equation}\label{tg:affine}
\begin{gathered}
 \dot G_{<3}+D^TG_{<3}=-dK_1^2B_1-dK_2^2\chi^2B_2,\qquad
 C:=D_{21}-D_{12}=2\dot\alpha+\Omega_1,\\
 \dot C-G_{<3,1}=2\ddot\alpha-A_0\sin\alpha-dK_1^2\Omega_1.
\end{gathered}
\end{equation}
Thus the full equations with equal physical diffusion are satisfied by
the explicitly specified scalar force, vector force and pressure
\begin{equation}\label{tg:affine_forces}
\begin{aligned}
 f_\vartheta&=(-dK_1^2B_1-dK_2^2\chi^2B_2)\cdot x,\\
 f_u&=\tfrac12(2\ddot\alpha-A_0\sin\alpha-dK_1^2\Omega_1)Jx,\\
 p&=-\tfrac12x^T\operatorname{sym}
       (\dot D+D^2-e_2\otimes G_{<3})x.
\end{aligned}
\end{equation}
To verify the new temperature contribution, differentiate $B_2$
using $\dot P_2=-dK_2^2\chi^2P_2$ and
$\dot\zeta_2=-D^T\zeta_2$; this gives
$\dot B_2+D^TB_2=-dK_2^2\chi^2B_2$.
The base-plus-first contribution is the exact identity proved in
Section~\ref{sr:section}, so their sum proves the first equation
in \eqref{tg:affine}. The curl identity is unchanged because
$\Omega_2=0$ and $B_{2,1}=0$, as the displayed vertical phase shows.
The trace of $D$ is zero, so $D^2$ is scalar. Substitution of the
quadratic pressure cancels the symmetric part of
$\dot D+D^2-e_2\otimes G_{<3}$; its skew part is
$(\dot C-G_{<3,1})J/2$. This proves the entire vector equation.
Both affine Laplacians vanish by second differentiation.
These forces are the affine-core parent forces. The compact released
profiles add their exact profile, envelope and interaction terms;
the present identities do not discard those spatial residuals or
claim an unforced third-layer superposition.
